% ===========================================================================
%  附录：补充分析与讨论
%  format2026.doc 第四条规定正文不超过 30 页、附录页数不限。正文篇幅已达上限，
%  故把稳健性、灵敏度、成因分解与适用边界这类"支撑性"讨论移入本附录。所有数据、
%  公式与结论一条未删；正文对应位置保留了结论摘要（含全部关键数字）并指向本节。
% ===========================================================================
\newpage
\section{补充分析与讨论}
\label{app:extra}

本附录给出基准对照、灵敏度检验、成因分解及关键模型设定的补充结果。

\subsection{问题一：基准对照与灵敏度分析}

\subsubsection{基准对照：储能与光伏的价值}

为量化储能设备与光伏发电的经济价值，构造"光伏供电、缺口直接向外网购买、
电池保持不动"的基准方案：$g_t^{\mathrm{base}} = \max\{\ell_t - v_t,\,0\}$，
$c_t = d_t = 0$，$E_t \equiv 6000$，多余光伏弃用。该方案满足同样的日循环
约束，因而是原问题的一个可行解，可用于检验优化结果并衡量储能的经济贡献。

\begin{table}[htbp]
  \centering
  \caption{不同方案的经济性对照}
  \label{tab:p1-benchmark}
  \begin{tabular}{lrrr}
    \toprule
    方案 & 全天购电量/kWh & 全天购电费/元 & 弃光电量/kWh \\
    \midrule
    无光伏无储能 & 111024.81 & 88319.36 & --- \\
    无储能       & 61789.94  & 48052.05 & 6247.96 \\
    本模型（含储能） & \textbf{59482.70} & \textbf{35126.95} & \textbf{0.00} \\
    \bottomrule
  \end{tabular}
\end{table}

由表~\ref{tab:p1-benchmark} 可见，配置储能设备后，全天购电费由 $48052.05$~元
降至 $35126.95$~元，节省 $12925.10$~元，降幅为 $26.90\%$；较
“无光伏无储能”方案节省 $53192.41$~元（$60.23\%$）。含储能方案的购电量由
$61789.94$~kWh 降至 $59482.70$~kWh，这是因为储能将无储能方案中的
$6247.96$~kWh 光伏盈余转移到无光照时段使用，从而提高光伏本地消纳率并降低
电网购电量。优化费用 $35126.95$~元低于基准费用
$48052.05$~元，与上述可行性论证一致。

\subsubsection{灵敏度分析}

题目只给出一个"充放电效率为 $90\%$"的参数，其在充电与放电两侧的分配方式
存在解释空间：可以是单程效率（$\eta_c=\eta_d=0.9$，往返 $0.81$），也可以是
往返效率（对称折算为 $\eta_c=\eta_d=\sqrt{0.9}\approx0.9487$，往返 $0.90$）。
此外，储能的日循环水平、充放电决策的时间粒度也均可能影响结论。为此设计
六组对照口径：

\begin{table}[htbp]
  \centering
  \caption{不同建模口径下的求解结果对照}
  \label{tab:p1-sensitivity}
  \small
  \begin{tabular}{lrrr}
    \toprule
    口径 & 全天购电量/kWh & 全天购电费/元 & 弃光/kWh \\
    \midrule
    \textbf{主模型}（往返 $0.81$，$E_0=6000$，$10$ min）
      & \textbf{59482.70} & \textbf{35126.95} & 0.00 \\
    放电无损（$\eta_c=0.90$，$\eta_d=1.00$，往返 $0.90$）
      & 57580.98 & 33416.53 & 0.00 \\
    对称折算（$\eta_c=\eta_d=\sqrt{0.90}$，往返 $0.90$）
      & 57526.24 & 33801.50 & 0.00 \\
    $E_0$ 自由（仅要求 $E_0=E_{144}$，收敛至 $8550$）
      & 59482.70 & 35118.60 & 0.00 \\
    粗化为 $1$ 小时粒度
      & 59924.95 & 36771.29 & 768.44 \\
    粗化为 $4$ 小时粒度
      & 61302.10 & 43648.06 & 2293.67 \\
    \bottomrule
  \end{tabular}
  \par\vspace{2pt}
  \begin{minipage}{\textwidth}
    \footnotesize 注："往返"指 $\eta_c\eta_d$；除注明外均为
    $\eta_c=\eta_d=0.90$、$E_0=6000$~kWh、$10$ 分钟决策粒度。
  \end{minipage}
\end{table}

对表~\ref{tab:p1-sensitivity} 的分析可得以下结论：

\textbf{（1）效率参数的影响。}若将 $90\%$ 的效率解释为往返
效率并对称折算（$\eta_c=\eta_d=\sqrt{0.9}$），购电费为 $33801.50$~元，较主
模型降低 $3.77\%$；若极端地假设放电无损（$\eta_c=0.90$，$\eta_d=1.00$），
购电费为 $33416.53$~元，降低 $4.87\%$。两者构成单程/往返两种解释下的费用
区间 $[33416.53,\ 35126.95]$~元，跨度不足 $5\%$。原因是放电效率越低，
储能放电时需额外多消耗的储存电量越多，越直接地削弱峰谷套利收益。\textbf{两
种口径下"谷充峰放"的策略结构与零弃光结论均不变}，说明主模型对效率参数的
分配方式不敏感。

\textbf{（2）初始储电量的影响。}若解除 $E_0 = 6000$~kWh 的
固定约束、仅要求日循环 $E_0 = E_{144}$，模型自动收敛到 $E_0 = 8550$~kWh，
购电费仅下降 $8.35$~元（$0.024\%$）。这表明在当前数据和约束下，初始储电量
对最优费用的影响较小，采用题目给定的 $E_0 = 6000$~kWh 不改变主要结果。

\textbf{（3）决策粒度的影响。}当
充放电决策由 $10$ 分钟粗化为 $1$ 小时时，购电费上升 $4.68\%$ 并产生
$768.44$~kWh 弃光；粗化为 $4$ 小时时，购电费上升 $24.26\%$（达
$43648.06$~元），弃光量进一步增至 $2293.67$~kWh。原因是粗粒度下储能功率在
各时段内保持恒定，对电价与光伏出力变化的响应能力下降，峰谷调节和光伏消纳
受到影响。因此，$10$ 分钟决策粒度能更充分地利用时段信息。

\subsection{问题二：弃用电量的成因}

全年弃用电量达 $2463544.98$~kWh，约占同期光伏发电量
$19068890.0$~kWh 的 $12.9\%$。题目未对弃电设置惩罚或收益，因此弃电不产生
额外结算费用，但可以反映日前计划与实际运行之间的偏差。

按逐时段口径分解全年弃电，可以看到它来自\textbf{两个性质不同}的来源。
$71.8\%$ 的弃电（$1769824.2$~kWh）发生在实际光伏大于实际负载的时段，属于
光伏富余超出储能吸纳能力——储能受 $5000$~kW 功率上限与 $10800$~kWh 容量上限
双重约束，在光伏最集中的时段无法全额吸收；风险修正又使日前计划量系统性偏高，
与高于预期的实际光伏叠加后进一步加剧了这一部分。另 $28.2\%$ 的弃电
（$693720.7$~kWh）发生在\textbf{光伏并无富余}的时段，其中 $137052.5$~kWh
甚至出现在光伏出力为零的深夜与凌晨——这部分电量不可能来自光伏，只能来自
\textbf{已按计划电价付费、但因实际负载低于预期而用不掉的计划购电量}。按各
时段电价核算，这部分被弃掉的已购电量当初共付出 $466267$~元，约占全年计划
购电费的 $3.54\%$。因此，这部分弃电对应为降低紧急购电风险而增加的计划购电
支出。

计划购电量同时影响弃电和紧急购电。本文在固定
$\rho=1$、$\lambda=0$ 下重跑了 $\alpha\in\{0.50,0.70,0.80,0.90\}$ 四点：$\alpha$
由 $0.80$ 降至 $0.50$ 时，弃电量从 $2628.6$~MWh 降至 $1493.7$~MWh，但紧急
购电量由 $177.9$~MWh 增至 $636.1$~MWh，总费用从 $1403.8$ 万元
升至 $1508.8$ 万元；反向升至 $0.90$ 时，紧急购电量降至 $86.2$~MWh，弃电量
却升至 $3369.1$~MWh，总费用回升至 $1421.6$ 万元；$\alpha=0.70$ 为
$1410.8$ 万元。四点中 $\alpha=0.80$ 取到最小值且两侧点均高于它，
表明风险余量需要在弃电与紧急购电之间权衡。（这四点均固定 $\rho=1$；
主策略因滚动重标定在 $8$ 个窗口选到不同参数，全年 $1393.97$ 万元还略低于该
网格最优点。）逐点记录见支撑材料。

图~\ref{fig:p2-daily} 给出紧急购电量的季节分布与逐日费用构成。紧急购电费用
在时间上较为集中。
$334$ 天中有 $143$ 天发生紧急购电，但其中位日仅 $317.1$~kWh；而最严重的
$6$ 月 $1$ 日单日紧急购电 $10545.11$~kWh、费用 $56753.15$~元，\textbf{一天就
占去全年紧急购电费的 $7.5\%$}，其后紧接的 $6$ 月 $2$ 日（$7244.83$~kWh、
$42449.26$~元）与 $12$ 月 $2$ 日、$7$ 月 $6$ 日同样显著高出常态。
这说明全年紧急购电费主要由少数连续高负荷时段构成。$6$ 月初的连日高负荷与
光伏出力波动同时出现，
使储能在数日内持续被压至下限，缺口只能在 $5$ 倍电价下补足。

这一分布表明，额外储能容量或预测资源可优先用于高缺口日期。

\subsection{问题二：模型设定与运行含义}

问题二以连续储能状态连接相邻运行日，以过去 $35$ 天的实际总费用滚动标定风险
余量与储备线。该设定把跨日储能价值通过连续状态和后续决策纳入全年仿真；在给定的
$5000$~kW 功率上限下，策略较无储能基准节省 $24.73\%$ 的费用。

\subsection{问题三：运行设定与结算规则}

每个六小时调整区间仅在区间起点提交一次调整，因此每个时段只相对初始计划
$g_t^0$ 结算一次。每次求解只使用当时已发布的信息，故 $14{,}977{,}809$~元的
对照费用与 $13{,}418{,}626$~元的主策略费用均对应各自可得的信息集。参数以
2025 年 1 月实际总费用标定，并在评价期固定；主口径下调减部分退回原价并收取
$50\%$ 违约费，不退款口径的独立结果列于表~\ref{tab:p3-billing}。

\subsection{问题四：预测与结算设定}
\label{app:p4-limit}

附件 4 与附件 1 的年平均相同，逐时段均值最大偏差仅
$5.15\times10^{-5}$~元/kWh；两者差值的星期结构解释 $41.7\%$ 的方差。模型据此
将日内形状、星期特征和价格加权误差分位数结合为价格预测与风险修正。第
\ref{sec:p4-risk} 节的正协方差结果表明，价格加权分位数能够将高价时段的缺口风险
直接纳入购电计划。主结果按退款口径结算：取消部分退回原价，并按交易时刻价格的
$50\%$ 支付违约费；不退款口径的独立时序仿真见附录第 \ref{app:p4-refund} 节。
